% !TEX root = ../main.tex
\section{Ejercicio propuesto 1: biblioteca de primalidad}

\subsection{Interfaz y criterio de decisión}
La biblioteca \path{python/primality.py} reúne las siete funciones pedidas. Su contrato distingue una prueba concluyente de composición de una aceptación provisional:
\begin{table}[H]\centering\small
\caption{Interfaz pública organizada por función}\label{tab:interfaz}
\begin{tabularx}{\textwidth}{@{}p{.16\textwidth}>{\raggedright\arraybackslash}p{.42\textwidth}X@{}}\toprule
Etapa & Funciones & Resultado principal \\ \midrule
Prueba exacta & \texttt{trial\_division} & Decisión determinista hasta $\lfloor\sqrt n\rfloor$; $n<2$ no es primo. \\
Rondas & \texttt{fermat\_round}\newline\texttt{decompose}\newline\texttt{miller\_rabin\_round} & Base válida $2\le a\le n-2$; $n-1=2^sd$. La ronda de Miller--Rabin devuelve decisión y detalle: $n,a,s,d$, residuos y contadores. \\
Decisión & \texttt{is\_probable\_prime} & Precribado seguido de rondas; aceptar significa \emph{probable primo}. \\
Generación & \texttt{generate\_candidate}, \texttt{generate\_probable\_prime} & Entero impar de longitud exacta; muestreo hasta superar el precribado y las rondas. \\ \bottomrule
\end{tabularx}\end{table}

Se comprueban tipos, bases y descomposiciones antes de operar. El precribado usa primos hasta $997$: certifica al candidato que coincide con uno de ellos y descarta sus múltiplos. La exponenciación binaria registra los cuadrados \emph{efectivamente realizados}; no atribuye a \texttt{pow()} un conteo interno.

\subsection{Casos y trazas: qué demuestra cada salida}
La Tabla \ref{tab:casos_propuesto1} procede de \path{python/results/casos_conocidos.csv}. ``Sí'' en Fermat o Miller--Rabin indica que \emph{pasa la base 2}; para $n<5$ la ronda no aplica. Los cuatro compuestos que pasan Fermat en la tabla son rechazados por Miller--Rabin con esa misma base.
\begin{table}[H]\centering\small
\caption{Casos especiales y pseudoprimos: resultado de las tres pruebas}
\label{tab:casos_propuesto1}
\begin{tabular}{rlccc}\toprule
$n$ & Clasificación & División & Fermat (2) & MR (2) \\
\midrule
0 & ni primo ni compuesto & no & -- & -- \\
1 & ni primo ni compuesto & no & -- & -- \\
2 & primo & sí & -- & -- \\
3 & primo & sí & -- & -- \\
25 & compuesto & no & no & no \\
341 & compuesto & no & sí & no \\
561 & compuesto & no & sí & no \\
1105 & compuesto & no & sí & no \\
1729 & compuesto & no & sí & no \\
\bottomrule\end{tabular}\end{table}


\begin{table}[H]\centering\small
\caption{Trazas verificables de Miller--Rabin con base $2$}
\label{tab:trazas_propuesto1}
\begin{tabular}{rrrrll}\toprule
$n$ & $a$ & $s$ & $d$ & Residuos & Decisión \\
\midrule
221 & 2 & 2 & 55 & [128, 30] & testigo \\
341 & 2 & 2 & 85 & [32, 1] & testigo \\
561 & 2 & 4 & 35 & [263, 166, 67, 1] & testigo \\
\bottomrule\end{tabular}\end{table}


En las trazas de la Tabla \ref{tab:trazas_propuesto1}, generadas desde \path{python/results/trazas_miller_rabin.csv}, $221$ no alcanza $-1$; $341$ y $561$ alcanzan $1$ desde una raíz no trivial. Las tres rondas aportan \emph{testigos de composición}.

\subsection{Validación y alcance}
Las \textbf{32 pruebas} de \path{python/test_primality.py} cubren extremos, cuadrados de primos y bases inválidas; comparan $10\,001$ enteros con una criba independiente y $1000$ aleatorios menores que $10^6$ con división de prueba. Además, \texttt{classify\_bases} separa testigos y mentirosos fuertes para todas las bases válidas de los compuestos impares $5\le n<200$ y de $221$, $341$ y $561$. Los candidatos de $16$ a $1024$ bits conservan longitud e imparidad.

\texttt{primality\_report} conserva bases, rondas y operaciones para auditar la decisión. Un testigo certifica composición; si un compuesto impar supera $r$ bases \emph{uniformes e independientes}, la probabilidad de aceptarlo no excede $4^{-r}$. La semilla fija solo reproduce ensayos: sin \texttt{rng} se usa \texttt{secrets.SystemRandom}, que sí es la fuente prevista para generación segura.
